\begin{abstract}
Em um avião monomotor, a falha do motor em voo elimina toda a propulsão.
Comparamos a confiabilidade de três tecnologias, representadas pelo Lycoming O-540 (pistão), pelo PT6A-34 (turboélice) e pelo CFM56-7B (turbofan), com diagramas de blocos de confiabilidade que incluem pares redundantes e falhas de causa comum.
A taxa total de cada motor vem da literatura, e a parcela de cada subsistema foi estimada em 1109 achados de componente de 2899 falhas de motor investigadas pelo NTSB entre 2008 e 2026.
Em um voo de 1~h, o motor a pistão falha com probabilidade de $1{,}24\times10^{-4}$, 83 vezes a do turboélice e 62 vezes a do turbofan.
A separação resiste a toda a faixa publicada.
No pistão, cilindros, válvulas, pistões e virabrequim respondem por 60\% das falhas, enquanto os magnetos duplicados ficam com 0,6\%; nas turbinas, compressores e turbinas lideram.
Os pares redundantes reduzem em cerca de 10 vezes a chance de perder a função, ganho limitado pela fração de causa comum $\beta$.
Dois motores a pistão só superam um turboélice se menos de 1,2\% das falhas forem comuns aos dois motores; contra o turbofan, o limite é 1,6\%.
No NTSB, 43\% das falhas de motor registradas em multimotores a pistão tiveram vítima fatal, contra 11\% nos monomotores a pistão e 3\% nos turboélices monomotores.
Um teste caso-controle com o tempo desde a revisão de 903 motores a pistão dá forma de Weibull 0,99, compatível com a taxa constante.
\end{abstract}

\begin{IEEEkeywords}
confiabilidade, motores aeronáuticos, diagrama de blocos de confiabilidade, falha de causa comum, desligamento em voo
\end{IEEEkeywords}
